The Voting Rights Act of 1965, particularly Section 2 as amended in 1982, prohibits electoral systems that dilute minority voting strength. Courts have interpreted this to require creation of majority-minority (MM) districts where geographically possible \citep{thornburg1986, bartlett2009}. However, achieving VRA compliance through principled, algorithmic methods that avoid partisan gerrymandering remains an open challenge.

Traditional redistricting approaches face a fundamental tension: VRA compliance requires concentrating minority populations into specific districts (non-uniform distribution), while principled redistricting emphasizes compactness and contiguity (favoring uniform distribution). Recent work on algorithmic redistricting using graph partitioning methods \citep{dellaLibera2026recursive, duchin2018, deford2021} focuses primarily on compactness and population balance, with limited attention to VRA constraints.

\subsection{Research Questions}

This paper investigates three fundamental questions:

\begin{enumerate}
\item \textbf{Feasibility}: Can principled, multi-constraint optimization methods achieve VRA compliance without resorting to partisan map-drawing?

\item \textbf{Geographic Constraints}: What demographic and geographic factors determine whether a state can achieve VRA targets through algorithmic methods?

\item \textbf{Methodological Comparison}: How do different optimization approaches (n-way partitioning, adaptive recursive bisection, constraint relaxation) perform for VRA compliance?
\end{enumerate}

\subsection{Key Contributions}

We make three primary contributions:

\textbf{(1) Comprehensive VRA Testing}: First systematic evaluation of multi-constraint optimization for VRA compliance across multiple covered states, using METIS graph partitioning with 2D vertex weights (population + minority voting-age population).

\textbf{(2) Critical Threshold Discovery}: We identify a critical threshold around 42\% state-wide minority population, above which VRA compliance is achievable with compact districts, below which it becomes infeasible without relaxing compactness.

\textbf{(3) Geographic Dominance}: We demonstrate that geographic clustering of minority populations is the primary determinant of VRA compliance feasibility, not algorithmic sophistication. Advanced methods (n-way partitioning, adaptive bisection) improve concentration by 3-7 percentage points but cannot overcome fundamental geographic constraints.

\subsection{Research Program Context}

This paper extends the foundational recursive bisection method presented in~\cite{della-libera-recursive-bisection}, which establishes algorithmic objectivity through structural immunity to partisan manipulation, and builds on the edge-weighted approach in~\cite{della-libera-edge-weighted}, which optimizes geometric compactness. While those baseline methods achieve population balance and compactness, they do not explicitly address Voting Rights Act requirements for majority-minority district creation. This work adds demographic awareness to the partitioning framework, investigating whether principled algorithmic methods can satisfy VRA obligations while maintaining the impossibility defense against gerrymandering. Subsequent papers in this research program investigate the demographic thresholds that determine VRA feasibility and the compactness tradeoffs inherent in creating majority-minority districts.

\subsection{Implementation and Results Note}

Results labeled ``V4 production'' in this paper reflect the adaptive formula achieving full legal targets in 2/5 covered states (Alabama and Georgia); 4/5 states achieve targets under per-state-optimal fixed configurations in the ablation study. Post-\emph{Callais} (2026), the VRASection algorithm (\texttt{--structure ratio-optimal-vra}) implemented in the \texttt{redist} Rust binary provides the recommended production implementation, superseding the Python-based workflow originally used to generate these results.

\subsection{Overview}

Section 2 reviews VRA requirements and prior algorithmic redistricting work. Section 3 describes our multi-constraint optimization methodology and four tested approaches. Section 4 presents results across five VRA-covered states. Section 5 discusses implications for redistricting practice and policy. Section 6 concludes with recommendations for future work.
